% document formatting
\documentclass[10pt]{article}
\usepackage[utf8]{inputenc}
\usepackage[left=1in,right=1in,top=1in,bottom=1in]{geometry}
\usepackage[T1]{fontenc}
\usepackage{xcolor}

% math symbols, etc.
\usepackage{amsmath, amsfonts, amssymb, amsthm}

% lists
\usepackage{enumerate}

% images
\usepackage{graphicx} % for images
\usepackage{tikz}

% code blocks
% \usepackage{minted, listings} 
\usepackage[colorlinks=true, urlcolor=blue]{hyperref}


% verbatim greek
\usepackage{alphabeta}

\graphicspath{{./assets/images/Week 4}}

\title{02-719 Week 4 \\ \large{Genomics and Epigenetics of the Brain}}
 
\author{Aidan Jan}

\date{\today}

\begin{document}
\maketitle

\section*{Diversity of Neuron Subtypes}
Neuroscience in 2024:
\begin{itemize}
    \item 15$+$ types of cortical interneurons with distinct morphologies, firing patterns, and connectivity
    \item 49 cortical cell populations with distinct transciptomes
    \item Just beginning to understnad the diversity of neurons on the epigenetic levevl and we are seeing vastly different epigenetic landscapes
\end{itemize}

\subsection*{Major Cell Types of the Brain}
\begin{itemize}
	\item Oligodendrocytes
	\begin{itemize}
	    \item Form the myelin sheath around axons, allowing for fast signal propagation
    \end{itemize}
	\item Microglia
	\begin{itemize}
	    \item Immune cells that detect damaged or unhealthy neurons and protect against foreign bacteria and viruses
    \end{itemize}
	\item Astrocytes
	\begin{itemize}
	    \item Provide structure to hold neurons in place
	    \item Supply nutrients
	    \item Digest debris from dead cells
    \end{itemize}
	\item Neurons
	\begin{itemize}
	    \item Facilitate information transfer through propagation of electrical signal
    \end{itemize}
\end{itemize}
\begin{center} 
	\includegraphics*[width=0.8\textwidth]{W4_1.png} 
\end{center}

\subsubsection*{Types of Neurons}
Neurons can be divided into two groups:
\begin{itemize}
	\item \textbf{Glutamatergic Pyramidal Neurons:} excitatory neurons.
	\begin{itemize}
	    \item Use glutamate as a chemical messenger to increase brain activity.  Receiving cells are more like to fire an electral signal
	    \item Drive most rapid signaling, thinking, learning, and memory.
	    \item Hyperactive glutamatergic neurons leads to conditions like epilepsy.
    \end{itemize}
	\item \textbf{GABAergic Interneurons:} inhibitory neurons
	\begin{itemize}
	    \item Use gamma-aminobutyric acid (GABA) as a chemical messenger to calm brain activity.  Receiving cells are less likely to fire an electrical signal
	    \item Prevent runaway seizures and balance brain signals.
    \end{itemize}
\end{itemize}
Each group contains more different types of neurons.
\begin{center} 
	\includegraphics*[width=0.9\textwidth]{W4_2.png} 
\end{center}

\subsection*{Cellular Diversity in the Cortex}
Excitatory Neurons: ``Principal Cells''
\begin{itemize}
	\item \textbf{Intracortical:}
	\begin{itemize}
	    \item Predominantly in Layer 2 / 3
	    \item Projections stay within the cortex
	    \item Can be \textit{associative} (project to the same hemisphere, area, or column that the cell body is in) or \textit{commissural} (project to the opposite hemisphere)
    \end{itemize}
	\item \textbf{Corticofugal:}
	\begin{itemize}
	    \item Mainly in Layer 5 / 6
	    \item Project to brain regions outside of the cortex
	    \item Can be \textit{corticalthalamic} (project to the thalamus, important for processing sensory information) or \textit{subcerebral} (project to the brainstem or spinal cord)
    \end{itemize}
\end{itemize}
Molecular interrogation reveals additional heterogeneity within subtypes of excitatory neurons.  Subtypes can be defined by characteristic expression levels of a set of genes

\subsection*{Different Views of Diversity}
\begin{center} 
	\includegraphics*[width=\textwidth]{W4_3.png} 
\end{center}

\section*{Issues with Bulk Tissues}
\begin{center} 
	\includegraphics*[width=0.6\textwidth]{W4_4.png} 
\end{center}
Signals from rare cell types can often be missed entirely.

\section*{Searching for the Genetic Basis of Disease}
\begin{itemize}
	\item We can run GWAS (genome-wide association study) to correlate SNPs with certain diseases.
	\item With enough reads, we can determine which peaks (SNPs) represent mutations that correlate with the presence of diseases
	\item Cell type-specific peaks overlap key disease SNPs, which might lead to problems.  Is a mutation a cause for a disease, or is it normal for a different cell type?
\end{itemize}

\section*{Single Cell Experiments}
To do sequencing studies, the vast majority of the data comes from single cell experiments.

\subsection*{Tissue Source}
\begin{itemize}
	\item \textbf{Fresh tissue:}
	\begin{itemize}
	    \item Best cell viable - best quality RNA
	    \item Requires direct access to tissue, lower throughput
    \end{itemize}
	\item \textbf{Young (embryonic) tissue:}
	\begin{itemize}
        \item Less myelinated tissue to dissociate
        \item Age is experimental confounder
    \end{itemize}
	\item \textbf{Cryo-preserved frozen:}
	\begin{itemize}
        \item Better cell viability, higher throughput
        \item Freeze-thaw effects in sample
    \end{itemize}
	\item \textbf{Flash frozen}
	\begin{itemize}
        \item Human brain repositories usually flash freeze tissue
        \item Worst cell viability
    \end{itemize}
\end{itemize}

\subsection*{Enzymatic Dissociation of Single Cells}
\begin{enumerate}
	\item Tissue is dissected and cut into thin slides.
	\item Tissue is rinsed with PBS (Phosphate Buffered Saline), then minced.
	\item Digestive enzymes are added (which might include collagenase, hyaluronidase, DNAse, Papain, etc.)
	\item Mixture is incubated to ensure tissue is digested
	\item Sample is inspected and filtered (so only nuclei are left)
	\item Nuclei can by input through a machine for cell counting and sequencing
\end{enumerate}

\subsection*{Differences Between Single Cell and Single Nuclei}
As it turns out, the amount of RNA from cells is more than the RNA in nuclei
\begin{itemize}
	\item This is because RNA is present in cytoplasm
\end{itemize}
However, the number of cell types detected from cells is equal to types detected from nuclei
\begin{itemize}
	\item Essentially, a cell's DNA sequence (and epigenetic markers) can be used to determine cell type
\end{itemize}
Both preparations can detect similar cell types.

\subsection*{Fresh Nuclei from Fresh Brain Tissue}
\begin{itemize}
	\item Less size bias in nuclei vs. cells
	\item Pestle mechanically break up extracellular matrix
	\item Iso/hypotonic buffer solution maintain organelle membranes (e.g., nuclei, mitochondria, etc.)
\end{itemize}
Nuclei are then purified using a density gradient.
\begin{itemize}
	\item Use solutions of 40\% iodixanol, 30\% iodixanol, and 25\% iodixanol.  Higher percentages are denser than lower percentages, and will therefore sink.
	\item Nuclei, being more dense, will sink to between the 30\% and 40\% layers.  Meanwhile, all the junk will float at the 25\% layer.
	\item A needle can be used to extract the nuclei.
\end{itemize}


\end{document}